\documentclass[11pt]{article}
\usepackage{amsmath,amssymb,amsthm}
\usepackage{geometry}
\usepackage{booktabs}
\usepackage{graphicx}
\usepackage{hyperref}
\usepackage{listings}
\usepackage{xcolor}
\usepackage{algorithm}
\usepackage{algpseudocode}

\geometry{letterpaper, margin=1in}

\lstset{
  basicstyle=\ttfamily\small,
  keywordstyle=\color{blue},
  commentstyle=\color{gray},
  stringstyle=\color{red},
  breaklines=true,
  columns=fullflexible,
  keepspaces=true,
  language=C++
}

\newtheorem{definition}{Definition}
\newtheorem{theorem}{Theorem}
\newtheorem{lemma}{Lemma}
\newtheorem{proposition}{Proposition}

\title{\textbf{Formation Engine Methodology} \\
\Large Section 10b — Defect Microstate Classification:\\
Third Regime for Occupancy and Disorder}
\author{Formation Engine Development Team}
\date{Version 0.1 — Unified Clustering Architecture}

\begin{document}
\maketitle

\section{Introduction: Beyond Perfect Crystals}
\label{sec:defect-intro}

Sections~10--13 established crystal structure generation, supercell construction,
and perfect lattice relaxation. Real materials, however, rarely occupy a single
ideal configuration. Experimental specimens exhibit:
\begin{itemize}
  \item \textbf{Vacancies and interstitials} (point defects)
  \item \textbf{Substitutional disorder} (dopants, solid solutions)
  \item \textbf{Antisite defects} (wrong atom on wrong sublattice)
  \item \textbf{Charge compensation patterns} (oxygen vacancies $\leftrightarrow$
        cation oxidation states)
  \item \textbf{Supercell-scale ordering motifs} (short-range order, clustering)
\end{itemize}

These features define a \textbf{continuum of microstates} within a parent phase,
not discrete polymorphs. This section introduces the \textbf{defect (microstate)
regime}---a third classification paradigm complementing polymorphs
(\S10: same chemistry, different structure) and isomorphs (\S10: same structure,
different chemistry) with a unified deterministic pipeline.

\section{Regime Definition and Invariants}

\subsection{Mathematical Formulation}

A structure instance is represented as:
\begin{equation}
\label{eq:struct-rep}
\mathcal{S} = (\mathbf{L}, \{\mathbf{r}_i\}, \{s_i\})
\end{equation}
where:
\begin{itemize}
  \item $\mathbf{L}$ is the lattice matrix ($3 \times 3$, columns are
        $\{\mathbf{a}, \mathbf{b}, \mathbf{c}\}$)
  \item $\mathbf{r}_i$ are fractional coordinates of atom $i$
  \item $s_i$ is the species label (atomic number or type identifier)
\end{itemize}

\begin{definition}[Defect Microstate Regime]
\label{def:defect-regime}
The defect microstate regime classifies structures by:
\begin{align}
\textbf{Invariant:} \quad & \text{parent prototype } \Pi \text{ (topology/space group)} \\
\textbf{Variable:} \quad & \text{occupancy configuration } \Omega \text{ (site decoration + vacancies)} \nonumber \\
& \text{local distortion field } \Delta \text{ (geometry perturbations)} \nonumber
\end{align}

Formally, let $\Pi(\mathcal{S})$ denote the prototype assignment (via
topology fingerprint). The defect class of $\mathcal{S}$ is:
\begin{equation}
\mathrm{DefectClass}(\mathcal{S}) = \mathrm{Cluster}_{\epsilon}
\big(\mathrm{FP}_{\mathrm{def}}(\mathcal{S}) \;\big|\; \Pi(\mathcal{S}) = \Pi_0\big)
\end{equation}
where clustering is performed \emph{within} a fixed parent prototype $\Pi_0$.
\end{definition}

\subsection{Regime Comparison}

Table~\ref{tab:regime-comparison} contrasts the three classification regimes:

\begin{table}[h]
\centering
\caption{Unified Classification Regimes}
\label{tab:regime-comparison}
\begin{tabular}{lccc}
\toprule
\textbf{Regime} & \textbf{Invariant} & \textbf{Variable} & \textbf{Output} \\
\midrule
Polymorph & Composition & Structure & $\alpha, \beta, \gamma$ phases \\
Isomorph & Prototype & Chemistry & Decoration variants \\
Defect Microstate & Parent phase & Occupancy + disorder & $\omega_1, \omega_2, \ldots$ \\
\bottomrule
\end{tabular}
\end{table}

\paragraph{Key insight:} All three regimes use the \emph{same clustering engine}
with different fingerprints. This unification enables:
\begin{itemize}
  \item One canonicalization routine
  \item One neighbor-building algorithm
  \item One deterministic clustering implementation
  \item Two fingerprint constructors: $\mathrm{FP}_{\mathrm{proto}}$ (structure)
        and $\mathrm{FP}_{\mathrm{deco}}$ (chemistry/defects)
\end{itemize}

\section{Deterministic Pipeline}

\subsection{Core Algorithm}

The defect classification pipeline (Algorithm~\ref{alg:defect-pipeline}) extends
the crystal supercell workflow (\S10) with microstate-aware fingerprints:

\begin{algorithm}[H]
\caption{Defect Microstate Classification}
\label{alg:defect-pipeline}
\begin{algorithmic}[1]
\Require Set of relaxed structures $\{\mathcal{S}_1, \ldots, \mathcal{S}_N\}$
\Ensure Defect class assignments $\{\omega_1, \ldots, \omega_N\}$
\State \textbf{Phase 1: Parent Prototype Assignment}
\For{each $\mathcal{S}_i$}
  \State $\mathcal{S}_i \gets \mathrm{canonicalize}(\mathcal{S}_i)$
        \Comment{Wrap coords, order atoms}
  \State $\mathbf{G}_i \gets \mathrm{build\_neighbors}(\mathcal{S}_i)$
        \Comment{Fixed cutoff rule}
  \State $\mathrm{FP}_{\mathrm{proto}}^{(i)} \gets
         \mathrm{proto\_fingerprint}(\mathbf{G}_i, \mathcal{S}_i)$
        \Comment{Topology hash + CN}
\EndFor
\State $\{\Pi_1, \ldots, \Pi_M\} \gets
       \mathrm{cluster}(\{\mathrm{FP}_{\mathrm{proto}}^{(i)}\}, \epsilon_{\mathrm{proto}})$
       \Comment{Prototype IDs}

\State
\State \textbf{Phase 2: Defect Microstate Clustering (within each $\Pi_k$)}
\For{each prototype $\Pi_k$}
  \State $\mathcal{I}_k \gets \{i : \Pi(\mathcal{S}_i) = \Pi_k\}$
        \Comment{Structures in parent $\Pi_k$}
  \For{each $i \in \mathcal{I}_k$}
    \State $\mathrm{FP}_{\mathrm{def}}^{(i)} \gets
           \mathrm{defect\_fingerprint}(\mathcal{S}_i, \Pi_k)$
           \Comment{Eq.~\eqref{eq:defect-fp}}
  \EndFor
  \State $\{\omega_{k,1}, \omega_{k,2}, \ldots\} \gets
         \mathrm{cluster}(\{\mathrm{FP}_{\mathrm{def}}^{(i)}\}_{i \in \mathcal{I}_k},
         \epsilon_{\mathrm{def}})$
\EndFor
\State \Return $\{(\Pi(\mathcal{S}_i), \omega(\mathcal{S}_i)) : i = 1, \ldots, N\}$
\end{algorithmic}
\end{algorithm}

\subsection{Fingerprint Construction}

\subsubsection{Prototype Fingerprint (Chemistry-Agnostic)}

The prototype fingerprint captures \emph{topology} and \emph{geometry} independent
of specific species:

\begin{equation}
\label{eq:proto-fp}
\mathrm{FP}_{\mathrm{proto}}(\mathcal{S}) =
\Big(
\mathrm{WL\_hash}(\mathbf{G}),\;
\mathbf{h}_{\mathrm{CN}},\;
\mathbf{h}_{\mathrm{RDF}}
\Big)
\end{equation}
where:
\begin{itemize}
  \item $\mathrm{WL\_hash}(\mathbf{G})$ is the Weisfeiler-Lehman graph hash
        (canonical topology signature)
  \item $\mathbf{h}_{\mathrm{CN}}$ is the coordination number histogram
        (per site class, if known)
  \item $\mathbf{h}_{\mathrm{RDF}}$ is a normalized radial distribution function
        (optional, for geometric refinement)
\end{itemize}

All components are computed deterministically from $\mathbf{G}$ (the neighbor
graph at fixed cutoff) and canonicalized coordinates.

\subsubsection{Defect Fingerprint (Chemistry + Occupancy)}

The defect fingerprint captures \emph{local disorder} within a fixed parent:

\begin{equation}
\label{eq:defect-fp}
\mathrm{FP}_{\mathrm{def}}(\mathcal{S}) =
\Big(
\mathbf{o},\;
\mathbf{v},\;
\mathbf{m},\;
\mathbf{d}
\Big)
\end{equation}
where:
\begin{align}
\mathbf{o} &= \text{site-class occupancy vector}
  && \text{(counts or fractions per Wyckoff position)} \nonumber \\
\mathbf{v} &= \text{vacancy/interstitial vector}
  && \text{(deviation from ideal site count)} \nonumber \\
\mathbf{m} &= \text{substitution/antisite signature}
  && \text{(sublattice decoration pattern)} \nonumber \\
\mathbf{d} &= \text{local distortion signature}
  && \text{(CN deviation + bond-length residual histograms)} \nonumber
\end{align}

\paragraph{Component definitions:}
\begin{enumerate}
  \item \textbf{Occupancy vector $\mathbf{o}$:}
    For each site class $\alpha \in \{\alpha_1, \ldots, \alpha_K\}$ (determined
    from parent prototype via symmetry or neighbor topology):
    \begin{equation}
    o_\alpha = \frac{N_\alpha}{\langle N_\alpha \rangle_{\text{ideal}}}
    \end{equation}
    where $N_\alpha$ is the actual count and $\langle N_\alpha \rangle_{\text{ideal}}$
    is the expected count from the parent prototype.

  \item \textbf{Vacancy vector $\mathbf{v}$:}
    \begin{equation}
    v_\alpha = 1 - o_\alpha
    \qquad \text{(negative indicates interstitial region)}
    \end{equation}

  \item \textbf{Substitution signature $\mathbf{m}$:}
    Species composition per site class, expressed as a histogram over atomic
    numbers (or type IDs). For site class $\alpha$:
    \begin{equation}
    m_{\alpha,Z} = \frac{\#\{\text{atoms of species } Z \text{ on sites of class } \alpha\}}
                        {N_\alpha}
    \end{equation}
    Antisite defects (A on B sublattice) manifest as nonzero $m_{\beta,A}$ where
    $\beta$ is expected to host species B.

  \item \textbf{Distortion signature $\mathbf{d}$:}
    For each atom $i$, compute:
    \begin{align}
    \mathrm{CN}_i &= \text{coordination number from neighbor graph} \\
    \langle r_i \rangle &= \text{mean nearest-neighbor distance}
    \end{align}
    Then build histograms:
    \begin{equation}
    \mathbf{d} = \big(\mathbf{h}_{\mathrm{CN}\text{-dev}}, \mathbf{h}_{r\text{-dev}}\big)
    \end{equation}
    where:
    \begin{align}
    \mathbf{h}_{\mathrm{CN}\text{-dev}} &: \text{histogram of }
      (\mathrm{CN}_i - \langle \mathrm{CN} \rangle_{\text{ideal}}) \\
    \mathbf{h}_{r\text{-dev}} &: \text{histogram of }
      (\langle r_i \rangle - \langle r \rangle_{\text{ideal}})
    \end{align}
\end{enumerate}

\subsection{Distance Metric and Clustering}

Define a weighted Euclidean distance over defect fingerprints:
\begin{equation}
\label{eq:defect-dist}
D_{\mathrm{def}}(\mathcal{S}_a, \mathcal{S}_b) =
w_o \|\mathbf{o}_a - \mathbf{o}_b\| +
w_v \|\mathbf{v}_a - \mathbf{v}_b\| +
w_m \,\delta(\mathbf{m}_a, \mathbf{m}_b) +
w_d \,\mathrm{dist}(\mathbf{d}_a, \mathbf{d}_b)
\end{equation}
where:
\begin{itemize}
  \item $\|\cdot\|$ is the $\ell^2$ norm for continuous vectors
  \item $\delta(\mathbf{m}_a, \mathbf{m}_b)$ is a discrete substitution distance
        (e.g., sum of absolute differences of histograms)
  \item $\mathrm{dist}(\mathbf{d}_a, \mathbf{d}_b)$ is the Wasserstein distance
        or $\chi^2$ statistic between histograms
  \item $(w_o, w_v, w_m, w_d)$ are fixed weights (default: uniform)
\end{itemize}

Clustering proceeds deterministically via threshold assignment:
\begin{equation}
\mathrm{Cluster}_{\epsilon}(\{\mathrm{FP}_1, \ldots, \mathrm{FP}_N\}) =
\{\omega_1, \ldots, \omega_N\}
\end{equation}
where:
\begin{enumerate}
  \item Process structures in canonical order (stable hash of $\mathrm{FP}_{\mathrm{proto}}$)
  \item For each structure $i$:
    \begin{itemize}
      \item Compute $D_{\mathrm{def}}(\mathrm{FP}_i, \mathrm{centroid}_k)$ for
            all existing clusters $k$
      \item If $\min_k D_{\mathrm{def}}(\mathrm{FP}_i, \mathrm{centroid}_k) < \epsilon$,
            assign to that cluster
      \item Otherwise, create new cluster with ID $\mathrm{hash}(\mathrm{FP}_i)$
    \end{itemize}
  \item Update centroids as running means
\end{enumerate}

This avoids stochasticity from seeding or iteration order. The cluster ID is a
deterministic hash, not a sequential integer.

\section{Labeling and Statistics}

\subsection{Symbolic Labels (Optional)}

Defect class labels are \emph{derived properties}, not inputs. If symbolic labels
(e.g., $\omega_1, \omega_2, \omega_3$) are desired, they are assigned
deterministically by:
\begin{equation}
\text{Rank clusters by:} \quad
\begin{cases}
\text{Mean energy per formula unit (if available)} \\
\text{Population (number of structures in cluster)} \\
\text{Lexicographic hash of centroid fingerprint}
\end{cases}
\end{equation}

Labels are then assigned $\omega_1, \omega_2, \ldots$ in rank order. This ensures
reproducibility: the same dataset always yields the same labels.

\subsection{Per-Class Statistics}

For each defect class $\omega_k$ within parent $\Pi$, compute:
\begin{align}
\langle E \rangle_{\omega_k} &= \frac{1}{N_k} \sum_{i \in \omega_k} E_i
  && \text{(mean energy per formula unit)} \\
\sigma_E^2(\omega_k) &= \frac{1}{N_k} \sum_{i \in \omega_k} (E_i - \langle E \rangle)^2
  && \text{(energy variance)} \\
f_{\omega_k} &= \frac{N_k}{N_{\text{total}}}
  && \text{(population fraction within } \Pi \text{)} \\
\langle \mathbf{o} \rangle_{\omega_k} &= \frac{1}{N_k} \sum_{i \in \omega_k} \mathbf{o}_i
  && \text{(mean occupancy vector)} \\
\langle \mathbf{v} \rangle_{\omega_k} &= \frac{1}{N_k} \sum_{i \in \omega_k} \mathbf{v}_i
  && \text{(mean vacancy fraction)}
\end{align}

These statistics enable:
\begin{itemize}
  \item \textbf{Thermodynamic interpretation:} Classes with lower $\langle E \rangle$
        and higher $f$ are thermodynamically favored under the simulation conditions
  \item \textbf{Defect concentration analysis:} $\langle \mathbf{v} \rangle$ quantifies
        typical vacancy/interstitial densities
  \item \textbf{Stability screening:} Classes with low $\sigma_E^2$ are
        structurally consistent; high $\sigma_E^2$ suggests heterogeneous microstates
\end{itemize}

\section{Unified Implementation}

\subsection{Shared Core Functions}

All three regimes (polymorphs, isomorphs, defects) share a common implementation:

\begin{lstlisting}[caption={Unified Clustering Core (C++20)}, label=lst:unified-core]
namespace atomistic {
namespace classify {

// Canonicalize structure: wrap coords to [0,1), order atoms
Structure canonicalize(const Structure& s);

// Build neighbor graph from fixed cutoff rule
NeighborGraph build_neighbors(const Structure& s, double cutoff);

// Fingerprints
struct ProtoFingerprint {
    uint64_t topology_hash;   // Weisfeiler-Lehman hash
    std::vector<int> CN_hist; // Coordination number histogram
    std::vector<double> RDF;  // Radial distribution function
};

struct DefectFingerprint {
    std::vector<double> occupancy;       // o vector
    std::vector<double> vacancy;         // v vector
    std::map<int, std::vector<double>> substitution;  // m (Z -> histogram)
    std::vector<double> CN_deviation;    // d: CN histogram
    std::vector<double> bond_deviation;  // d: bond length histogram
};

ProtoFingerprint compute_proto_fp(const NeighborGraph& g, const Structure& s);
DefectFingerprint compute_defect_fp(const Structure& s, const Prototype& parent);

// Distance functions
double distance(const ProtoFingerprint& a, const ProtoFingerprint& b);
double distance(const DefectFingerprint& a, const DefectFingerprint& b,
                const std::array<double, 4>& weights = {1,1,1,1});

// Clustering (deterministic threshold assignment)
template<typename FP>
std::vector<uint64_t> cluster(const std::vector<FP>& fingerprints,
                               double epsilon);

} // namespace classify
} // namespace atomistic
\end{lstlisting}

\subsection{Module A: Polymorph Detection}

\begin{lstlisting}[caption={Polymorph Classifier (Same Chemistry $\to$ Structure Phases)}]
std::map<std::string, std::vector<uint64_t>>
classify_polymorphs(const std::vector<Structure>& structures) {
    // Group by reduced formula
    auto grouped = group_by_formula(structures);
    
    std::map<std::string, std::vector<uint64_t>> results;
    
    for (auto& [formula, structs] : grouped) {
        std::vector<ProtoFingerprint> fps;
        for (auto& s : structs) {
            auto canon = canonicalize(s);
            auto graph = build_neighbors(canon, FIXED_CUTOFF);
            fps.push_back(compute_proto_fp(graph, canon));
        }
        
        // Cluster on prototype fingerprint (ignore chemistry)
        auto clusters = cluster(fps, PROTO_EPSILON);
        
        // Rank by energy → assign α, β, γ labels
        auto labels = assign_greek_labels(clusters, structs);
        
        results[formula] = labels;
    }
    
    return results;
}
\end{lstlisting}

\subsection{Module B: Isomorph Detection}

\begin{lstlisting}[caption={Isomorph Classifier (Same Structure $\to$ Chemistry Variants)}]
std::map<uint64_t, std::vector<uint64_t>>
classify_isomorphs(const std::vector<Structure>& structures) {
    std::map<uint64_t, std::vector<uint64_t>> results;
    
    // Phase 1: Assign prototype IDs
    std::vector<ProtoFingerprint> proto_fps;
    for (auto& s : structures) {
        auto canon = canonicalize(s);
        auto graph = build_neighbors(canon, FIXED_CUTOFF);
        proto_fps.push_back(compute_proto_fp(graph, canon));
    }
    
    auto proto_clusters = cluster(proto_fps, PROTO_EPSILON);
    
    // Phase 2: Within each prototype, cluster by chemistry (decoration)
    for (uint64_t proto_id : unique(proto_clusters)) {
        auto indices = where(proto_clusters, proto_id);
        
        std::vector<DefectFingerprint> deco_fps;
        for (size_t i : indices) {
            deco_fps.push_back(compute_defect_fp(structures[i], proto_id));
        }
        
        auto variant_clusters = cluster(deco_fps, DECO_EPSILON);
        results[proto_id] = variant_clusters;
    }
    
    return results;
}
\end{lstlisting}

\subsection{Module C: Defect Microstate Detection}

\begin{lstlisting}[caption={Defect Classifier (Same Parent $\to$ Occupancy/Disorder Classes)}]
std::map<uint64_t, std::vector<uint64_t>>
classify_defects(const std::vector<Structure>& structures) {
    // Phase 1: Assign parent prototype IDs (same as isomorph Phase 1)
    auto proto_clusters = assign_prototypes(structures, PROTO_EPSILON);
    
    std::map<uint64_t, std::vector<uint64_t>> results;
    
    // Phase 2: Within each parent, cluster by defect fingerprint
    for (uint64_t parent_id : unique(proto_clusters)) {
        auto indices = where(proto_clusters, parent_id);
        
        std::vector<DefectFingerprint> defect_fps;
        for (size_t i : indices) {
            defect_fps.push_back(compute_defect_fp(structures[i], parent_id));
        }
        
        // Weights: emphasize vacancy + substitution over pure distortion
        std::array<double, 4> weights = {1.0, 1.5, 2.0, 0.5};  // o, v, m, d
        
        auto defect_clusters = cluster(defect_fps, DEFECT_EPSILON, weights);
        
        // Rank by population or mean energy → assign ω₁, ω₂, ω₃ labels
        auto labels = assign_omega_labels(defect_clusters, structures[indices]);
        
        results[parent_id] = labels;
    }
    
    return results;
}
\end{lstlisting}

\section{Example: MgO with Vacancy Disorder}

\subsection{System Setup}

Consider a simulation generating 500 relaxed MgO structures at 1200~K with:
\begin{itemize}
  \item 90\% stoichiometric MgO (rocksalt structure, $Fm\overline{3}m$)
  \item 8\% structures with $\sim$1\% Mg vacancies
  \item 2\% structures with O vacancies (charge-compensated by Mg$^+$)
\end{itemize}

\subsection{Classification Pipeline}

\paragraph{Step 1: Prototype Assignment.}
All 500 structures cluster to \textbf{one prototype ID} (rocksalt FCC sublattice
topology with octahedral/tetrahedral voids).

\paragraph{Step 2: Defect Fingerprinting.}
For each structure, compute $\mathrm{FP}_{\mathrm{def}}$:
\begin{itemize}
  \item $\mathbf{o} = (o_{\mathrm{Mg}}, o_{\mathrm{O}})$ where $o_{\mathrm{Mg}} \approx 0.99$
        for vacancy-containing structures
  \item $\mathbf{v} = (1 - o_{\mathrm{Mg}}, 1 - o_{\mathrm{O}})$ captures vacancy fractions
  \item $\mathbf{m}$ captures whether Mg or O sites host vacancies (antisite-like)
  \item $\mathbf{d}$ shows distortion histograms (CN deviations near vacancies)
\end{itemize}

\paragraph{Step 3: Clustering.}
With $\epsilon_{\mathrm{def}} = 0.05$, clustering yields:
\begin{itemize}
  \item $\omega_1$: 450 structures (stoichiometric, CN=6 everywhere)
  \item $\omega_2$: 40 structures (Mg vacancies, $\langle v_{\mathrm{Mg}} \rangle \approx 0.01$)
  \item $\omega_3$: 10 structures (O vacancies, $\langle v_{\mathrm{O}} \rangle \approx 0.01$)
\end{itemize}

\paragraph{Step 4: Statistics.}
\begin{align}
\langle E \rangle_{\omega_1} &= -10.3 \;\text{eV/fu} && \text{(baseline stoichiometric)} \\
\langle E \rangle_{\omega_2} &= -10.1 \;\text{eV/fu} && \text{(+0.2~eV penalty for Mg vacancy)} \\
\langle E \rangle_{\omega_3} &= -9.8 \;\text{eV/fu} && \text{(+0.5~eV penalty for O vacancy)}
\end{align}

Population fractions:
\begin{equation}
f_{\omega_1} = 0.90, \quad f_{\omega_2} = 0.08, \quad f_{\omega_3} = 0.02
\end{equation}

This automatically recovers the defect chemistry from simulation output, with no
manual classification or pre-specified labels.

\section{Validation and Robustness}

\subsection{Reproducibility}

\paragraph{Determinism:}
All fingerprints, distances, and cluster assignments are deterministic. Running
the pipeline twice on the same dataset produces identical $\omega$ labels, even
if structure order changes (because canonicalization sorts atoms and cluster IDs
are hashes).

\paragraph{Threshold sensitivity:}
Vary $\epsilon_{\mathrm{def}}$ from $0.01$ to $0.10$ and verify:
\begin{itemize}
  \item $\epsilon$ too small $\to$ over-clustering (every structure is its own class)
  \item $\epsilon$ too large $\to$ under-clustering (all structures merge to one class)
  \item Optimal $\epsilon$ yields stable cluster count across small perturbations
\end{itemize}

Use the elbow method or silhouette score to auto-tune $\epsilon$ (while maintaining
determinism by fixing random seeds to zero or using grid search with fixed order).

\subsection{Physical Consistency}

\paragraph{Energy ordering:}
For a correct force field and sufficient sampling:
\begin{equation}
\langle E \rangle_{\omega_1} < \langle E \rangle_{\omega_2} < \langle E \rangle_{\omega_3}
\quad \Leftrightarrow \quad
f_{\omega_1} > f_{\omega_2} > f_{\omega_3}
\end{equation}
This is a \textbf{necessary condition} for thermodynamic validity: lower-energy
defect classes should appear more frequently in equilibrium sampling.

\paragraph{Vacancy concentration bounds:}
For point-defect-dominated regimes:
\begin{equation}
\langle v_\alpha \rangle_{\omega_k} < 0.10
\quad \text{(typically)}
\end{equation}
Classes with $\langle v \rangle > 0.20$ likely represent amorphous regions or
phase transitions, not simple defect microstates.

\subsection{Comparison to Manual Classification}

Run the pipeline on a known dataset with expert-assigned labels (e.g., from DFT
calculations with annotated defect types). Compute:
\begin{equation}
\mathrm{Purity} = \frac{1}{K} \sum_{k=1}^K \frac{\max_{\ell} |\omega_k \cap L_\ell|}{|\omega_k|}
\end{equation}
where $L_\ell$ are the expert labels. Purity $> 0.90$ indicates strong agreement.

\section{Limitations and Honest Assessment}

\paragraph{Site-class determination.}
The pipeline assumes site classes (Wyckoff positions or sublattices) are either:
\begin{enumerate}
  \item Known from the parent prototype (via space group)
  \item Inferred from neighbor topology (via graph partitioning)
\end{enumerate}
For highly disordered or amorphous regions, site-class assignment may be ambiguous.
In such cases, fall back to global composition histograms (no per-site resolution).

\paragraph{Large supercells.}
For supercells with $N > 1000$ atoms, fingerprint computation scales as
$\mathcal{O}(N^2)$ for neighbor graphs. Optimization strategies:
\begin{itemize}
  \item Use cell lists or Verlet lists for $\mathcal{O}(N)$ neighbor building
  \item Subsample atoms for defect fingerprints (e.g., every 10th atom for $\mathbf{d}$)
  \item Cache topology hashes to avoid recomputation
\end{itemize}

\paragraph{Force field dependence.}
Energy-based ranking ($\langle E \rangle_{\omega_k}$) depends on force field accuracy.
For classical potentials (LJ, Buckingham), expect qualitative correctness
(relative energy ordering) but not quantitative precision (absolute formation energies).

\paragraph{Non-equilibrium sampling.}
If structures are quenched from high temperature or kinetically trapped, population
fractions $f_{\omega_k}$ do \emph{not} reflect equilibrium thermodynamics. The
classification is still valid (microstates are correctly grouped), but
energy-vs-population correlations require equilibration.

\section{Implementation Status}

\begin{itemize}
  \item \textbf{Core fingerprints:} Implemented in \texttt{atomistic/classify}
        (prototype + defect fingerprints)
  \item \textbf{Clustering:} Deterministic threshold assignment in
        \texttt{atomistic/classify/cluster.cpp}
  \item \textbf{Polymorph module:} Tested on 50-structure H$_2$O dataset
        (reproduces known ice phases)
  \item \textbf{Isomorph module:} Prototype, awaiting multi-chemistry validation
  \item \textbf{Defect module:} Interface designed, awaiting integration with
        \texttt{atomistic/crystal} supercell builder
\end{itemize}

\section{Conclusion: One Clustering Engine, Three Front-Ends}

This section extends the crystal classification framework (\S10) to handle
realistic disorder. The key contributions are:

\begin{enumerate}
  \item \textbf{Unified clustering architecture:} Three regimes (polymorphs,
        isomorphs, defects) share the same deterministic pipeline with different
        fingerprints
  \item \textbf{Defect fingerprint:} Four-component vector
        $(o, v, m, d)$ capturing occupancy, vacancies, substitutions, and
        distortion
  \item \textbf{Parent-aware clustering:} Defect classification is performed
        \emph{within} a fixed prototype, avoiding conflation of different
        structural basins
  \item \textbf{Deterministic labeling:} Symbolic labels ($\omega_1, \omega_2, \ldots$)
        are derived from energy/population ranking, not manual annotation
  \item \textbf{Validation protocol:} Reproducibility, energy ordering, and
        expert-label purity metrics ensure physical consistency
\end{enumerate}

\paragraph{Transition to production:}
With polymorphs, isomorphs, and defects unified under one clustering engine,
the next step is \textbf{automated phase diagram construction} (\S12):
given a composition range and temperature/pressure grid, iteratively generate
structures, classify them, and map the resulting phase space.

\end{document}
